\chapter{Critical review and limitations}
\label{ch:critique}

The chapter retrospectives give my view of each step as it happened. This chapter
consolidates what the work can and cannot be said to establish, and what would
have to be done to remove each limitation. The order is from the limitations that
bound the conclusions most tightly to those that bound them least.

Table~\ref{tab:recap} first puts every iteration on the one axis that survives a
change of corpus, the gap to the persistence baseline of its own test users. It
is the shortest complete account of the work in this document, and what follows
is a judgement about how much weight each of its rows can carry.

% T14 -- Every headline run of the internship on one axis, gap by gap.
\begin{table}[t]
\centering\footnotesize
\caption{Every iteration, on one axis}
\label{tab:recap}
\begin{tabular}{@{}c@{\hspace{8pt}}l@{\hspace{8pt}}l@{\hspace{8pt}}l@{\hspace{8pt}}r@{\hspace{8pt}}r@{\hspace{9pt}}r@{}}
\toprule
\textbf{It.} & \textbf{Corpus} & \textbf{Lever of that iteration} & \textbf{Protocol} & \textbf{Base.} & \textbf{Model} & \textbf{Gap} \\
\midrule
1  & \sraw     & pipeline built    & last 20\,\% & 69.2\,\% & 16.9\,\% & $-52.3$ \\
2  & \sraw     & harness fixes     & last 20\,\% & 69.2\,\% & 64.4\,\% & $-4.8$  \\
3  & \srawbig  & ten times more users & last 20\,\% & 81.6\,\% & 76.0\,\% & $-5.6$  \\
4  & \sfilt    & filtering         & last 20\,\% & 79.0\,\% & 80.1\,\% & $\mathbf{+1.1}$ \\
5  & \smob     & a cleaner corpus  & last 20\,\% & 45.5\,\% & 44.6\,\% & $-0.9$  \\
6  & \smob     & hour tokens       & last 20\,\% & 45.5\,\% & 44.2\,\% & $-1.3$  \\
7  & \smob     & group token       & full seq.   & 43.1\,\% & 41.3\,\% & $-1.8$  \\
8  & \smob     & adapter routing   & full seq.   & 43.1\,\% & 46.8\,\% & $\mathbf{+3.7}$ \\
9  & \stwok    & 1{,}900 users     & last 20\,\% & 48.2\,\% & 52.4\,\% & $\mathbf{+4.2}$ \\
9  & \slev     & levelled records  & last 20\,\% & 76.5\,\% & 74.9\,\% & $-1.6$  \\
10 & \seleven  & 21{,}000 users    & last 20\,\% & 50.2\,\% & 54.5\,\% & $\mathbf{+4.3}$ \\
11 & \seleven  & analysis only     & \multicolumn{4}{@{}l@{}}{\textit{no new training: why the ladder climbs, then stops}} \\
\bottomrule
\end{tabular}

\figtext{%
The eleven iterations reduced to the only quantity that is comparable between
them, the gap in points between the model and the persistence baseline of its
own test users. Reading down the \textbf{Gap} column is reading the whole
internship: five iterations below the trivial rule, one accidental win at
iteration~4, and a real one from iteration~8 onward. Reading down the
\textbf{Base.} column shows why no other column may be compared vertically,
since the corpus, and with it the difficulty of the test users, changed four
times. What each iteration actually did is in
Appendix~\ref{app:timeline}.}

\vspace{4pt}
\begin{minipage}{\textwidth}\footnotesize
``Last 20\,\%'' scores the final fifth of each test sequence, ``full seq.''
every position of it, one at a time, which is strictly harder and lowers the
baseline on the very same users from 45.5\,\% to 43.1\,\%
(Chapter~\ref{ch:time}). Iteration~6 is the row whose single number understates
it: its targeted sampling won $+2.6\pt$ inside the two guessed windows and lost
more than that outside them. Iteration~9 appears twice because two independent
things were tried that week, more users and levelled records.
\end{minipage}
\end{table}


\section{What the work establishes, and with what confidence}

\paragraph{Solid.} A fine-tuned model beats the persistence baseline by
$+3.7\pt$ on the mobile-user corpus and by $+4.2$ to $+4.7\pt$ on the two larger
corpora, five independent training runs, two different held-out test sets,
each measured against its own baseline. The $+4.3\pt$ quoted in the abstract is
one of those measurements, the largest corpus at its largest tier, 54.5\,\%
against 50.2\,\%. The effect is larger than the difference two identical runs
would plausibly show, and it is reproduced across corpora. The mechanism is identified: the
gain is concentrated in the behavioural profiles where repetition fails, and it
ranks in exact inverse order of the baseline across all four of them.

\paragraph{Solid, and negative.} From 1{,}900 to 21{,}000 training users
the edge over the baseline does not grow. This is not an absence of evidence but
a measured null: four tiers, a paired bootstrap over 1{,}000 held-out users with
2{,}000 resamples, and every confidence interval containing zero, including
the one spanning the whole range, whose point estimate is $+0.04\pt$. Six
different metrics give the same verdict.

\paragraph{Well supported but not proven.} The explanation for \emph{why} the
curve climbs between 400 and 1{,}900 users and then stops, that the smallest
behavioural profile crosses roughly a hundred members exactly there, fits the
timing and is consistent with the ablation showing the model ignoring an
under-populated group token. It is one mechanism that fits, not a mechanism that
has been isolated.

\paragraph{Measured, but not a gain.} Levelling every user to 200 records does
what it claims (per-user weight inequality falls to zero) and reverses the sign
of the gap ($+3.7 \rightarrow -1.6\pt$), for a reason that is arithmetic rather
than modelling. The diagnosis behind it is sound and the remedy is incompatible
with this pipeline.

\section{Limitations that bound the conclusions}

\paragraph{What saturates is a shared regularity, not a personal one.} Adding
users buys the model what those users have in common, and the flatness of the
scaling curve says that this shared part is a small object, bought early and
then exhausted. Nothing in the curve bounds what a model could learn from a
single person's own habits, because no run here was given a person's history to
learn them from.

\paragraph{Corpus-to-corpus comparisons are deltas only.} The test population
changes whenever the corpus does, and with it the baseline, which is why the
$+3.7\pt$ of iteration 8 and the $+4.2\pt$ of iteration 9 are two separate
measurements of the same effect rather than a progression. The clean fix is cheap
and was not done: re-scoring the earlier checkpoint on the later test file is a
single evaluation pass, no retraining, and it would put every row of
Table~\ref{tab:recap} on one axis.

\paragraph{The scaling null holds at constant compute.} Every tier read 8{,}000
training examples regardless of how many were available, so at 21{,}000 users
14{,}522 of 22{,}522 windows were never read once. The defensible claim is
``with this training budget, 1{,}900 users are enough'', and the experiment that
would separate it from ``more users never help'', giving the largest tier the
compute its size deserves, was not run.

\paragraph{Each configuration was trained once.} The paired bootstrap covers
the luck of which test users are drawn. It does not cover the other source of
chance, that training has a random component of its own, so that the same
configuration run twice does not land on exactly the same score. Differences of
a few tenths of a point between tiers, and the $+0.8\pt$ measured on the
easiest behavioural profile, therefore sit on a floor nobody measured.
Repeating one tier from a different random start would establish that floor, at
a cost of about half an hour.

\paragraph{Part of the gain may be capacity rather than grouping.} Routing four
behavioural groups to four adapters multiplies the trainable parameters by four.
The control, a single adapter carrying as many trainable numbers as the four
together, same data, was not run.
The ablations of the failed run argue that the grouping carries most of the
effect, since the model extracted nothing from the group when it arrived as a
token, but that is an argument and not a measurement.

\paragraph{The mechanism is weakest where it would be worth most.} This is the
limitation I would put first if someone asked what the work is good for. The
model's edge is largest on users about whom little is known, $+4.6\pt$ with only
30\,\% of a day in context against $+4.3\pt$ with 80\,\%
(Chapter~\ref{ch:scale}), and that cold start is exactly the case an operator
cares about: predicting the rest of the day for a user who has just appeared is
worth more than confirming that a well-known one is still at home. But the group
a user belongs to is recovered from their context prefix with 92\,\% agreement at
the reference protocol and only 61\,\% and 44\,\% when half a day or 30\,\% of it
is available (Figure~\ref{fig:group-difficulty}(b)), so the very conditioning
that produces the edge decays as the history shortens. The two effects pull
against each other, and this report measures each separately without measuring
their combination. The defensible operating envelope for what is described here
is therefore a user with most of a day already recorded, and lifting that means
building a group detector that works on a short prefix, which is a smaller and
better-posed problem than the one I spent the internship on.

\paragraph{The metric deliberately penalises a non-movement.} The prediction
target is the raw \cid{cell\_id}, so a handset re-attaching to another cell of
the same antenna is scored as an error even though nobody moved, 12.4\,\% of
consecutive pairs are of that kind. This is a considered choice, explained in
Chapter~\ref{ch:user}: it is the question that was asked, and changing the target
midway would have made every earlier number incomparable. The consequence is
that a known, small share of the reported error is not a modelling failure.

\paragraph{One study area, one fortnight, one operator.} All results come from a
single Czech region over fifteen days in March 2014, with 369 cells over 113
sites. A denser urban area with thousands of cells, or a corpus in which
vocabulary coverage is not already complete, could behave differently,
particularly the saturation result, which partly reflects how small this world
is.

\paragraph{The day-of-week experiment did not finish.} A seven-day comparison,
one independent training per weekday, was prepared and launched; the session
crashed after five hours with none of the seven complete. No day-of-week effect
was measured, and the report claims none.

\section{What I would do next, in order}

\begin{enumerate}
  \item \textbf{Re-score the earlier checkpoint on the later test file.} One
        evaluation pass, no training. It is the cheapest possible improvement to
        the report's internal consistency and it removes the only caveat that
        affects the headline numbers.
  \item \textbf{Re-measure the switch-rate correlation on the final model.} Also
        one evaluation pass. The $r = -0.971$ of Chapter~\ref{ch:user} was
        measured on a model that was imitating the trivial rule, and a model that
        genuinely helps the volatile minority has to break the lockstep between a
        user's switch rate and that user's score, so the coefficient should come
        down in absolute value. It is the cheapest independent test of this
        report's main claim, taken on a thousand test users rather than on four
        profiles, and it would say whether the mechanism reaches individual users
        or only shows up once they are averaged into groups.
  \item \textbf{Repeat one tier from a different random start.} Half an hour,
        and it establishes the floor below which none of the small differences
        in this report should be read.
  \item \textbf{Extend the question from one day to several}, on a corpus whose
        identifiers are stable across days. This is the only item on the list
        that changes what the problem \emph{is} rather than how well it is
        measured, and it is the natural continuation of this work: predicting
        day $N{+}1$ from days $1 \ldots N$, and with it any notion of
        personalisation.
\end{enumerate}
